\documentclass[11pt,a4paper]{article}
\usepackage[margin=19mm,headheight=15pt]{geometry}
\usepackage{fontspec}
\setmainfont{Noto Serif}
\usepackage{amsmath,amssymb,mathtools}
\usepackage{xcolor}
\usepackage{fancyhdr}
\usepackage{enumitem}
\usepackage{microtype}
\setlength{\parindent}{0pt}
\setlength{\parskip}{6pt}
\definecolor{navy}{HTML}{183153}
\definecolor{lightgray}{HTML}{F3F5F7}
\pagestyle{fancy}\fancyhf{}
\lhead{\small ME 5309 -- Mid-Semester Examination}
\rhead{\small ME 5309}
\cfoot{\thepage}
\begin{document}
\begin{center}
{\Large\bfseries ME 5309 -- Mid-Semester Examination}\\[3pt]
{\large Answer Booklet}\\[8pt]
\end{center}
\textbf{Instructions:} Attempt any three questions. Each question carries 10 marks. Show working where appropriate.
\section*{Question 1 \hfill [10 marks]}
\subsection*{Q1(a) \hfill [4 marks]}

The pixel pitch is 0.005 mm, so f\_x = 16/0.005 = 3200 pixels.
Sensor width = 6.4 mm.
Belt field of view = 400x6.4/16 = 160 mm; spatial resolusion = 160/1280 = 0.125 mm/pixel.

\subsection*{Q1(b) \hfill [3 marks]}

The smallest resolvable force is 0.122 N.

\subsection*{Q1(c) \hfill [3 marks]}

\[\begin{aligned}F = 0.8187, G = 0.0906\end{aligned}\]

\section*{Question 2 \hfill [10 marks]}
\subsection*{Q2(a) \hfill [2 marks]}

Choose x₁ = θ and x₂ = θ̇.
Then ẋ₁ = x₂ and ẋ₂ = [τ - b x₂ - mgℓ sin(x₁)]/J.
At equilibrium:\[\begin{aligned}x_2 = 0 and τ_0 = mg\ell sin \theta_e\end{aligned}\]
Since |sin θ\_e| less than or equal to 1, an equilibrium exists when |τ₀| less than or equal to mgℓ = 2.4525 N m.\par\smallskip At rest $\dot\theta=0$ and $\ddot\theta=0$, so the damping term vanishes. Rearranging gives $\tau_0=mg\ell\sin\theta_e$. Here $mg\ell=1\times9.81\times0.25=2.4525$ N\,m. Since $-1\leq\sin\theta_e\leq1$, the required torque must lie between $-2.4525$ and $+2.4525$ N\,m.

\subsection*{Q2(b) \hfill [2 marks]}

The Jacobian has A₁₂ = 1 and A₂₂ = -b/J = -2.5.
The input matrix is B = [0, 50]ᵀ.

\subsection*{Q2(c) \hfill [4 marks]}

\vspace{2mm}

\subsection*{Q2(d) \hfill [2 marks]}

The gravity term changes the equilibrium, so the steady-state error is still nonzero.
The poles should change.

\section*{Question 3 \hfill [10 marks]}
\subsection*{Q3(a) \hfill [2 marks]}

First, numerical differentiation amplifies high-frequency encoder noise and quantisation because a small position error is divided by the short sample interval.
Second, differencing is sensitive to sampling/timing and can produce a noisy, jittery velocity signal; filtering reduces noise but adds lag.\par\smallskip In a finite difference, the change in position is divided by the sample interval. Encoder quantization error appears in the difference of two measurements, so it can be magnified by the division by $T_s$. The result may jump around even if the actual motion is smooth.

\subsection*{Q3(b) \hfill [2 marks]}

For H(s)=ωc/(s+ωc), ωc=2π(20)=125.664 rad/s and T=0.001 s.
Tustin gives H(z)=α(1+z⁻¹)/[(1+α)+(α-1)z⁻¹], α=ωcT/2=0.0628319.
Thus:\[\begin{aligned}y[k]=0.881765y[k-1]+0.0591175(x_k+x[k-1])\end{aligned}\]
DC gain = (2x0.0591175)/(1-0.881765) approximately 1.000.\par\smallskip $\omega_c=2\pi(20)=125.664$ rad/s. With $T=0.001$ s, $\alpha=\omega_cT/2=125.664(0.001)/2=0.062832$. The recurrence coefficients are approximately $0.881765$ on the previous output and $0.0591175$ on each input sample. Check: $2(0.0591175)/(1-0.881765)\approx1$, so the DC gain is close to one.

\subsection*{Q3(c) \hfill [4 marks]}

With L=[ℓ₁,ℓ₂]ᵀ and C=[1,0], A-LC=[[-ℓ₁,1], [-ℓ₂,-2]].
Its characteristic polynomial is det(sI-(A-LC)) = (s+ℓ₁)(s+2)+ℓ₂ = s²+(ℓ₁+2)s+(2ℓ₁+ℓ₂).
Desired poles at -10 give (s+10)²=s²+20s+100.
Hence ℓ₁+2=20 gives ℓ₁=18, and 2ℓ₁+ℓ₂=100 gives ℓ₂=64.\par\smallskip Expanding $(s+\ell_1)(s+2)+\ell_2$ gives $s^2+(\ell_1+2)s+(2\ell_1+\ell_2)$. The desired polynomial is $(s+10)^2=s^2+20s+100$. Matching the $s$ coefficient gives $\ell_1=18$; matching the constant term then gives $36+\ell_2=100$, hence $\ell_2=64$.

\subsection*{Q3(d) \hfill [2 marks]}

An observer uses the plant model and known control input to estimate unmeasured states such as velocity, and its poles can be designed.
A simple low-pass filter cannot exploit the state-space model in this way.
A filter may be safer when the model is poorly known or unreliable, because it relies less on model accuracy.\par\smallskip The observer uses the plant model and the measured output to reconstruct the unmeasured state. Its convergence depends on the model and the observer poles. A low-pass filter is simpler, but generally introduces phase lag and does not explicitly use the plant dynamics.

\section*{Question 4 \hfill [10 marks]}
\subsection*{Q4(a) \hfill [2 marks]}

Poles are -2±2j and Nff=8.

\subsection*{Q4(b) \hfill [2 marks]}

The actual feedforward contribution is 0.75x8=6.

\subsection*{Q4(c) \hfill [4 marks]}

I obtain s³+(a+bk₂)s²+(c+bk₁)s-bK\_I.
Desired polynomial is s³+8s²+24s+32, so k₂=6, k₁=18, K\_I=-32.

\subsection*{Q4(d) \hfill [2 marks]}

At steady state, a stable closed loop with integral action drives ξ̇=r-θ to zero; therefore r-θ=0 and the tracking error is zero, assuming no saturation or other implementation limits prevent convergence.
A practical failure mode is integrator windup during actuator saturation; the standard remedy is anti-windup, such as back-calculation or conditional integration.\par\smallskip With integral action, a constant tracking error keeps accumulating in the integral state, so the controller can adjust its output until the error goes to zero, assuming the closed loop is stable and the actuator is not constrained. Under saturation the integral state can keep growing; anti-windup limits or back-calculates that state.

\end{document}